\documentclass[11pt,letterpaper]{article}
\usepackage[margin=1in,headheight=15pt]{geometry}
\usepackage[T1]{fontenc}
\usepackage[utf8]{inputenc}
\usepackage{amsmath,amsthm}
\usepackage{newpxtext,newpxmath}
\usepackage{microtype}
\usepackage{mathtools}
\usepackage{booktabs,array,longtable}
\usepackage{xcolor}
\usepackage{enumitem}
\usepackage{fancyhdr}
\usepackage{aliascnt}
\usepackage{titlesec}
\usepackage[most]{tcolorbox}
\usepackage{listings}
\usepackage{hyperref}
\usepackage[nameinlink,noabbrev]{cleveref}
\definecolor{Olive}{HTML}{536443}
\definecolor{Ink}{HTML}{273128}
\definecolor{Soft}{HTML}{F3F4EC}
\definecolor{Rule}{HTML}{BDC6AC}
\definecolor{Muted}{HTML}{646B60}
\hypersetup{colorlinks=true,linkcolor=Olive,citecolor=Olive,urlcolor=Olive,
  pdftitle={Exact maximal order types of finite-alphabet transfinite words},
  pdfauthor={Research manuscript prepared with ChatGPT},
  pdfsubject={A proposed solution to a finite-alphabet tightness question in order theory}}
\titleformat{\section}{\Large\bfseries\color{Olive}}{\thesection}{0.7em}{}
\titleformat{\subsection}{\large\bfseries\color{Ink}}{\thesubsection}{0.7em}{}
\pagestyle{fancy}
\fancyhf{}
\fancyhead[L]{\small\color{Muted}Finite-alphabet transfinite words}
\fancyhead[R]{\small\color{Muted}Research manuscript}
\fancyfoot[C]{\small\thepage}
\renewcommand{\headrulewidth}{0.35pt}
\setlength{\parskip}{0.35em}
\setlength{\parindent}{1.2em}
\setlist{nosep,leftmargin=1.6em}
\setcounter{tocdepth}{1}
\allowdisplaybreaks[2]
\emergencystretch=2em
\newtheorem{theorem}{Theorem}[section]
\newaliascnt{lemma}{theorem}
\newtheorem{lemma}[lemma]{Lemma}
\aliascntresetthe{lemma}
\newaliascnt{proposition}{theorem}
\newtheorem{proposition}[proposition]{Proposition}
\aliascntresetthe{proposition}
\newaliascnt{corollary}{theorem}
\newtheorem{corollary}[corollary]{Corollary}
\aliascntresetthe{corollary}
\theoremstyle{definition}
\newaliascnt{definition}{theorem}
\newtheorem{definition}[definition]{Definition}
\aliascntresetthe{definition}
\newaliascnt{example}{theorem}
\newtheorem{example}[example]{Example}
\aliascntresetthe{example}
\theoremstyle{remark}
\newaliascnt{remark}{theorem}
\newtheorem{remark}[remark]{Remark}
\aliascntresetthe{remark}
\newcommand{\W}{\mathcal W}
\newcommand{\I}{\mathcal I}
\newcommand{\Lcal}{\mathcal L}
\newcommand{\len}[1]{\lvert #1\rvert}
\newcommand{\ev}{\operatorname{ev}}
\newcommand{\rankht}{\operatorname{ht}}
\newcommand{\otp}{\operatorname{o}}
\newcommand{\down}{\mathord{\downarrow}}
\newcommand{\up}{\mathord{\uparrow}}
\newcommand{\natprod}{\mathbin{\otimes}}
\newcommand{\natsum}{\mathbin{\oplus}}
\newcommand{\emb}{\preccurlyeq}
\newcommand{\nemb}{\npreccurlyeq}
\newcommand{\N}{\mathbb N}
\newtcolorbox{resultbox}{enhanced,breakable,colback=Soft,colframe=Rule,
  boxrule=0.5pt,arc=1.5mm,left=10pt,right=10pt,top=7pt,bottom=7pt}
\lstset{basicstyle=\ttfamily\small,breaklines=true,columns=fullflexible,
  backgroundcolor=\color{Soft},frame=single,rulecolor=\color{Rule},
  xleftmargin=0.5em,xrightmargin=0.5em,showstringspaces=false}

\begin{document}
\hypersetup{pageanchor=false}
\begin{titlepage}
\noindent{\small\scshape\color{Olive}Order theory \; / \; Ordinal combinatorics}
\vspace{0.9cm}

\noindent{\Huge\bfseries\color{Ink}Exact maximal order types\par
of finite-alphabet\par
transfinite words\par}
\vspace{0.45cm}
\noindent{\Large\color{Olive}A separator construction for all lengths below $\omega^k$}
\vspace{0.8cm}

\noindent Prepared for Vladimir Reshetnikov\par
\noindent September 19, 2026
\vspace{0.8cm}

\begin{abstract}
For a finite poset $P$, let $\W_k(P)$ be the quasi-order of $P$-labelled words
of ordinal length less than $\omega^k$, ordered by subsequence embedding.
Let $N_k(P)$ be the number of equivalence classes of nonempty indecomposable
words in this range. We give a proof of the formula
\[
  \otp(\W_k(P))=\omega^{\omega^{N_k(P)-1}}
  \qquad (|P|\ge2,\ 1\le k<\omega).
\]
The number $N_k(P)$ is finite and effectively computable from a canonical
poset of bounded-height periodic terms. In particular, if $J(P)$ counts all
order ideals, including the empty ideal, then
\[
  \otp(\W_2(P))=\omega^{\omega^{|P|+J(P)-2}}.
\]
For the two-element antichain this is $\omega^{\omega^4}$, the upper bound
whose tightness is left unverified in Altman's Example~3.15. For the two-element
chain it is $\omega^{\omega^3}$. The lower-bound argument does not assume
that concatenation preserves incomparability: it uses ordinal-length
protection to prevent absorption at separators. The package includes full
proofs, exact atom-poset computations, and reproducible symbolic embedding tests.
\end{abstract}

\vfill
\begin{resultbox}
\small
\textbf{Research status.} This is a proposed solution with a complete written
proof, not an independently refereed or machine-formalized result.
The source problem was checked against arXiv:2409.03199v2, the version listed
on September~19, 2026. The bounded-height tree/indecomposable correspondence
is established background; the claimed contribution here is attainment of
the finite-generator upper bound by protected separator constructions.
\end{resultbox}
\end{titlepage}
\hypersetup{pageanchor=true}
\pagenumbering{roman}
\tableofcontents
\clearpage
\pagenumbering{arabic}

\section{The problem and the claimed answer}\label{sec:intro}

The starting point is a very concrete question about two letters.
Write $A_2=\{a,b\}$ with $a$ and $b$ incomparable. A word is allowed to have
any ordinal length less than $\omega^2$: a finite number of infinite blocks,
possibly followed by a finite suffix. One word embeds in another when its
positions can be mapped increasingly to positions carrying the same letters.
How large can the order type of a well-ordered linear extension of this
embedding order be, after mutually embeddable words are identified?

Altman gives the upper bound $\omega^{\omega^4}$ and explicitly leaves its
attainment unverified in Example~3.15 of \cite{Altman}. In the same example,
he observes that the two-element chain has an improved upper bound
$\omega^{\omega^3}$, because two of the periodic building blocks coincide
up to embedding equivalence. The arXiv record lists version~2, revised
March~10, 2026; the example is on page~10 of that version. These are the
specific finite-alphabet questions addressed here.

There are two reasons to attack the binary case through a more general
statement. First, the mechanism responsible for an additional ordinal
exponent is not specific to either binary alphabet. Second, once the
correct lower-bound construction is found, the same construction can be
iterated at every finite ordinal level.

\subsection{Main theorem}

Precise definitions follow in \cref{sec:background,sec:atoms}. For now,
$\I_k(P)$ denotes the poset of equivalence classes of nonempty
indecomposable $P$-words of lengths
\[
  1,\ \omega,\ \omega^2,\ \ldots,\ \omega^{k-1},
\]
and $N_k(P)=|\I_k(P)|$. A word is \emph{indecomposable} here when it is
mutually embeddable with every nonempty final segment of itself.

\begin{theorem}[Finite-alphabet exactness]\label{thm:main}
Let $P$ be a finite poset with $n\ge2$ elements, and let $k\ge1$ be finite.
Then $\I_k(P)$ is finite and
\begin{equation}\label{eq:main}
  \boxed{\otp(\W_k(P))=\omega^{\omega^{N_k(P)-1}}.}
\end{equation}
Writing $J(R)$ for the number of order ideals of a finite poset $R$,
including the empty ideal, the numbers are computed by
\begin{equation}\label{eq:recurrence-intro}
  \I_0(P)=\varnothing,\qquad N_1(P)=n,\qquad
  N_{k+1}(P)=n+J(\I_k(P))-1.
\end{equation}
The recurrence requires the actual order on $\I_k(P)$, not just its cardinality.
\end{theorem}

The full proof is in \cref{sec:mainproof}. Two consequences can already be stated.

\begin{corollary}[The first transfinite level]\label{cor:first}
If $P$ is a finite poset of size $n\ge2$, then
\begin{equation}\label{eq:first}
  \otp(\W_2(P))=\omega^{\omega^{n+J(P)-2}}.
\end{equation}
In particular,
\[
  \otp(\W_2(A_2))=\omega^{\omega^4},\qquad
  \otp(\W_2(C_2))=\omega^{\omega^3},
\]
where $C_2$ is the two-element chain.
\end{corollary}

\begin{corollary}[Necessary exceptional cases]\label{cor:exception}
For the empty alphabet, $\otp(\W_k(\varnothing))=1$.
For the one-element alphabet,
\[
  \otp(\W_k(1))=\omega^k.
\]
A finite quasi-order is first replaced by its quotient under mutual comparison;
the number of classes, not the number of original labels, determines whether
these exceptions apply.
\end{corollary}

\begin{proof}
There is only the empty word over the empty alphabet. Over one letter,
a word is determined by its length up to equivalence, and embedding is
exactly ordinal comparison of lengths. The set of lengths less than
$\omega^k$ has order type $\omega^k$.
\end{proof}

\subsection{What is established background, and what is being added?}

The maximal-order-type calculus for well-quasi-orders is due to de Jongh
and Parikh and to subsequent work, particularly Schmidt
\cite{deJonghParikh,Schmidt}. Higman's finite-word construction and its exact
finite-alphabet order type are not new here \cite{Higman,Schmidt}.
Chopra and Pakhomov establish a correspondence between finite leaf-labelled
trees and indecomposable finite-image sequences below $\omega^\omega$
\cite[Theorem~4.26]{ChopraPakhomov}. The normal forms used below are a
bounded-height, antichain-reduced presentation of that existing structure.

The point requiring a new argument is the lower bound in
\cref{thm:main}. A finite list of indecomposable types gives an immediate
Higman upper bound. It does \emph{not} give a matching lower bound by simply
concatenating those types: for example,
\[
  a\,a^\omega\equiv a^\omega.
\]
The proof below constructs large induced suborders on which this absorption
is prevented by a final segment of sufficient ordinal length. Every new
indecomposable type can then serve as a separator and contribute one more
step in the finite exponent parameter.

The claim of novelty is limited to that exactness argument and its
consequences. The sources consulted did not supply this finite-level
attainment theorem; this is not an exhaustive determination of priority.
The theorem for finite alphabets is also not a claim to settle every
infinite-alphabet lower-bound question.

\section{Words, quotients, and ordinal accounting}\label{sec:background}

\subsection{Transfinite words and embedding}

\begin{definition}
A $P$-word of length $\alpha$ is a function $u\colon\alpha\to P$.
For words $u\colon\alpha\to P$ and $v\colon\beta\to P$, write $u\emb v$
if there is a strictly increasing map $f\colon\alpha\to\beta$ such that
$u(\xi)\le_P v(f(\xi))$ for every $\xi<\alpha$.
Set $u\equiv v$ when both $u\emb v$ and $v\emb u$.
\end{definition}

Concatenation is indexed by ordinal addition: if $|u|=\alpha$ and
$|v|=\beta$, then $|uv|=\alpha+\beta$. Finite and countable concatenations
preserve embedding in every component. In particular, replacing the
components of a concatenation by equivalent words produces an equivalent
word.

For a positive integer $k$,
\[
  \W_k(P)=\{u:|u|<\omega^k\},
\]
with the empty word included. Since $P$ is finite, every such word has finite
image. This is the object denoted $s^F_{\omega^k}(P)$ in the source problem.
All maximal order types below refer to the quotient by $\equiv$.

An increasing map between ordinals can exist only from a smaller ordinal
to a larger one. Consequently
\begin{equation}\label{eq:lengthmonotone}
  u\emb v\ \Longrightarrow\ |u|\le|v|,
  \qquad u\equiv v\ \Longrightarrow\ |u|=|v|.
\end{equation}
This elementary fact makes word length a well-defined invariant of an
equivalence class. It is ordinal length, not cardinality, that is relevant:
a bounded interval in $\omega^2$ may be countably infinite while still being
too short to contain an increasing copy of $\omega^2$.

\subsection{Maximal order type}

A quasi-order is a \emph{well-quasi-order} (wqo) if every infinite sequence
$x_0,x_1,\ldots$ has $i<j$ with $x_i\le x_j$. Its maximal order type
$\otp(X)$ is the largest order type of a linear extension of its quotient
poset. The existence of a largest type, rather than merely a supremum,
is part of the de Jongh--Parikh theory \cite{deJonghParikh}.
We use the following standard facts:
\begin{equation}\label{eq:wqocalculus}
\begin{aligned}
  \otp(X\sqcup Y)&=\otp(X)\natsum\otp(Y),\\
  \otp(X\times Y)&=\otp(X)\natprod\otp(Y),\\
  \otp(X)&=\sup_{x\in X}\bigl(\otp(X\setminus\up x)+1\bigr).
\end{aligned}
\end{equation}
Here $\sqcup$ is the disjoint union, the product is componentwise, and
$\oplus,\otimes$ are natural ordinal addition and multiplication.
An induced suborder or, more generally, the domain of an order-reflecting
map into a wqo has no larger maximal order type. A monotone surjection
from a wqo onto another quasi-order, also allowing surjectivity up to
equivalence, gives the reverse comparison of domain and image types.
These are standard tools in the same calculus \cite{deJonghParikh,Schmidt}.

For an integer $m\ge1$ put
\begin{equation}\label{eq:h}
  H_m=\omega^{\omega^{m-1}}.
\end{equation}
The finite-alphabet Higman formula is
\begin{equation}\label{eq:higman}
  \otp(R^*)=H_{|R|}\qquad(0<|R|<\omega),
\end{equation}
where $R^*$ is the finite-word order over the finite poset $R$.
For completeness, \cref{app:higman} derives the finite-alphabet formula
from the standard wqo calculus. The order relations among the letters do
not change \eqref{eq:higman}; they will matter for the transfinite atom count.

\subsection{The ordinal amplification identity}

The numerical calculation used repeatedly below is simple but important.
For $m,d\ge1$,
\begin{equation}\label{eq:power}
  H_m^{\otimes d}=\omega^{\omega^{m-1}\cdot d}.
\end{equation}
Indeed, multiplying monomials naturally adds their exponents naturally,
and the natural sum of $d$ copies of $\omega^{m-1}$ is
$\omega^{m-1}\cdot d$. The countable \emph{ordinary} ordinal sum satisfies
\begin{equation}\label{eq:amplify}
  \sum_{d=1}^{\infty}H_m^{\otimes d}
  =\omega^{\omega^{m-1}\cdot\omega}
  =\omega^{\omega^m}=H_{m+1}.
\end{equation}
The same expression is obtained by indexing $d=m'+1$ for $m'\ge0$.
No countable natural product is being invoked.

\begin{lemma}[Layered products]\label{lem:layers}
Suppose a wqo $Y$ contains induced copies $Y_d$ of $E^d$ for all $d\ge1$,
where $\otp(E)=H_m$. Suppose the copies are pairwise disjoint on the
quotient and
\[
  x\in Y_d,\ y\in Y_e,\ x\le y\quad\Longrightarrow\quad d\le e.
\]
Then $\otp(Y)\ge H_{m+1}$.
\end{lemma}

\begin{proof}
For each $d$, choose a maximal linear extension of $Y_d$. Order the layers
by increasing $d$ and use the chosen extension within each layer.
The hypothesis guarantees that this is a linear extension of their
\emph{induced union}. Its order type is the left side of
\eqref{eq:amplify}, by the product rule. The induced union is a suborder of
$Y$, so its type gives a lower bound for $\otp(Y)$.
There is no assertion that this particular linear extension must extend
to all of $Y$ with the same type.
\end{proof}

\section{Canonical indecomposables at finite height}\label{sec:atoms}

This section supplies a self-contained bounded-height normal form.
It is included to make the integer $N_k(P)$ effective and to identify
precisely which generators are used by the lower-bound proof. The general
tree/sequence correspondence is established in \cite{ChopraPakhomov}.

\subsection{Indecomposability and finite concatenations}

For a word $u\colon\alpha\to P$ and $\xi<\alpha$, its tail at $\xi$ is
the restriction to $[\xi,\alpha)$, reindexed by that interval's order type.

\begin{definition}
A nonempty word $u$ is indecomposable if it is equivalent to every one of
its nonempty tails.
\end{definition}

If $|u|<\omega^k$ and $u$ is indecomposable, then $|u|=\omega^r$ for some
$r<k$. To see this, write $|u|$ in Cantor normal form. Unless it is a single
power of $\omega$ with coefficient one, it has a nonempty tail of smaller
order type. That contradicts \eqref{eq:lengthmonotone}.
The property is invariant under equivalence: an embedding between two
words of the same indecomposable ordinal length is cofinal, so a tail in
one word receives a tail of the other.

\begin{lemma}[An indecomposable lands in one factor]\label{lem:onefactor}
If $z$ is indecomposable and
\[
  z\emb u_1u_2\cdots u_t
\]
for a finite list of words, then $z\emb u_i$ for some $i$.
\end{lemma}

\begin{proof}
Among the finitely many target factors met by the embedding, choose the
last one. From any source position whose image lies in that factor onward,
all source positions map into the same factor. Thus a nonempty tail of
$z$ embeds in it. That tail is equivalent to $z$.
\end{proof}

The lemma does not say that the original embedding places the entire
source inside one factor. It says that \emph{some} embedding does so.
This distinction will also matter for separators.

\subsection{Periodic terms and their order}

Start with a leaf $p$ of height zero for each $p\in P$.
Given a finite nonempty set $S$ of previously constructed terms, introduce
\[
  \Omega(S),\qquad
  \rankht(\Omega(S))=1+\max_{s\in S}\rankht(s).
\]
We will retain only sets $S$ that are antichains under the already defined
term order. Choose any enumeration $S=\{s_1,\ldots,s_t\}$ and define
\begin{equation}\label{eq:evaluation}
  \ev(p)=p,\qquad
  \ev(\Omega(S))=
  \bigl(\ev(s_1)\cdots\ev(s_t)\bigr)^\omega.
\end{equation}
The particular enumeration changes neither the embedding class nor the
comparison rule below. Every evaluation is indecomposable, and
\begin{equation}\label{eq:heightlength}
  |\ev(s)|=\omega^{\rankht(s)}.
\end{equation}
For a node, every tail contains all sufficiently late complete periods,
which proves indecomposability. The length calculation follows by taking
the largest height among the finitely many factors of a period and then
multiplying the period length by $\omega$ on the right.

\begin{proposition}[Comparison of periodic terms]\label{prop:comparison}
The following recursive rules agree exactly with semantic embedding:
\begin{align}
  p\le q &\quad\Longleftrightarrow\quad p\le_P q,
       \label{eq:leafleaf}\\
  p\le\Omega(T) &\quad\Longleftrightarrow\quad
       (\exists t\in T)\ p\le t,
       \label{eq:leafnode}\\
  \Omega(S)\le p &\quad\text{never},
       \label{eq:nodeleaf}\\
  \Omega(S)\le\Omega(T) &\quad\Longleftrightarrow\quad
       (\forall s\in S)(\exists t\in T)\ s\le t.
       \label{eq:nodenode}
\end{align}
In particular, comparison terminates by descending to lower-height terms.
\end{proposition}

\begin{proof}
The leaf rules follow directly from the presence or absence of a suitable
letter. A positive-height word cannot embed into a word of length one.

For necessity in \eqref{eq:nodenode}, fix one occurrence of a source factor
$\ev(s)$ inside a period of $\ev(\Omega(S))$. There is a later source
position, so the image of this factor is bounded in the target word.
A bounded initial segment of a countable repetition of a finite period
meets only finitely many complete or partial periods. Therefore
$\ev(s)$ embeds into a finite concatenation of target factors $\ev(t)$.
By \cref{lem:onefactor}, it embeds into one of them. This holds for each
$s\in S$.

Conversely, choose for each $s$ a suitable $t$ with $s\le t$.
Place successive source factors into suitably chosen target factors in
successive fresh target periods. Repeating this construction through
$\omega$ source periods gives an embedding. The argument with $S=T$
but different enumerations proves independence of the enumeration.
\end{proof}

For a finite set $S$ of terms, remove every element strictly below another
element of $S$, retaining the maximal elements. This changes neither side
of the Hoare comparison in \eqref{eq:nodenode}. Thus every node has an
antichain-reduced form. Moreover, if two finite antichains dominate one
another, then they are equal: for $s\in S$, choose $t\in T$ and $s'\in S$
with $s\le t\le s'$; the antichain property forces $s=s'$, and
antisymmetry forces $s=t$.

Induction on height now proves that the reduced terms form a genuine
poset, with no distinct but equivalent terms. Leaves and nodes cannot be
equivalent because their lengths differ.

\subsection{Completeness and finiteness}

\begin{theorem}[Finite decomposition and canonical atoms]\label{thm:normalform}
For every finite $P$ and integer $k\ge1$:
\begin{enumerate}
\item there are finitely many reduced terms of height less than $k$;
\item every $P$-word of length less than $\omega^k$ is equivalent to a
finite concatenation of evaluations of these terms;
\item every indecomposable word of length $\omega^r$, $r<k$, is equivalent
to exactly one reduced term of height $r$.
\end{enumerate}
Consequently those terms, with the order in \cref{prop:comparison}, are
precisely $\I_k(P)$.
\end{theorem}

\begin{proof}
For $k=1$ the words are finite and the only indecomposables are letters.
Assume the result through height $r-1$ and consider a word $u$ of length
$\omega^r$, where $r\ge1$.

Partition its domain into $\omega$ consecutive intervals of length
$\omega^{r-1}$. By induction, each resulting block is equivalent to a
finite nonempty list of lower-height atoms. Replacing all blocks by those
lists and flattening yields an $\omega$-indexed list of atoms of heights
at most $r-1$. There are only finitely many available atom types.

Let $S$ be the nonempty set of types occurring infinitely often in that
list. Delete a finite initial part containing all occurrences of every
type outside $S$. The remaining list is equivalent to a cyclic repetition
of $S$: map each of its successive entries into a fresh target period in
one direction, and choose successive sufficiently late occurrences of the
required entries in the other. Replacing $S$ by its maximal elements
gives a reduced node $\Omega(S)$.

Every original block has length $\omega^{r-1}$, so its finite factorization
must contain an atom of height $r-1$. Finitely many types are available;
hence some such type recurs infinitely often. Therefore the reduced set
$S$ still has maximal height $r-1$, and its node has height $r$.
We have shown that $u$ is equivalent to a finite lower-height prefix
followed by one atom of height $r$.

If $u$ itself is indecomposable, it is equivalent to that last atom.
One way to justify this without assuming a particular factorization of
$u$ is to use \cref{lem:onefactor}: the finite prefix factors have length
less than $\omega^r$, so $u$ cannot embed into them. It must embed into
the final atom, which already embeds into $u$.

For an arbitrary ordinal $\alpha<\omega^{r+1}$, split $\alpha$ into finitely
many blocks of length $\omega^r$ and one remainder of smaller length.
The preceding argument handles the blocks and induction handles the
remainder. This proves finite decomposition at the next level.
There are only finitely many nonempty antichains of the finite poset of
lower-height terms, so there are finitely many new terms. Uniqueness and
exact comparison follow from \cref{prop:comparison} and antichain reduction.
\end{proof}

\subsection{The exact atom-count recurrence}

Every order ideal of a finite poset is determined by its antichain of
maximal elements, and every antichain determines an ideal. Empty objects
correspond to one another. Thus $J(R)$ also counts the antichains of $R$.

All nonleaf terms of height at most $k$ are in bijection with the nonempty
antichains of $\I_k(P)$. This includes lower-height nodes already present:
a node with lower-height children is the same node, not a new copy.
Adding the $n$ leaves gives
\begin{equation}\label{eq:count}
  N_{k+1}(P)=n+J(\I_k(P))-1.
\end{equation}
Equivalently, the number of \emph{new} height-$k$ atoms is
$J(\I_k(P))-J(\I_{k-1}(P))$. This formulation avoids double-counting old nodes.

\begin{corollary}[The finite-generator upper bound]\label{cor:upper}
If $P$ is nonempty, then
\[
  \otp(\W_k(P))\le H_{N_k(P)}.
\]
More generally, if $B$ is any finite nonempty set of distinct indecomposable
word classes and $\Lcal(B)$ is its finite-concatenation language, then
$\otp(\Lcal(B))\le H_{|B|}$.
\end{corollary}

\begin{proof}
Concatenation gives a monotone map from $B^*$ onto $\Lcal(B)$ up to
equivalence. Apply the Higman formula and the monotone-surjection rule.
For $B=\I_k(P)$, \cref{thm:normalform} identifies the image with $\W_k(P)$.
The same argument proves that all these languages are wqo.
\end{proof}

\section{Cofinal blocks and protected final segments}\label{sec:blocks}

Fix throughout this section
\[
  \lambda=\omega^r,\qquad 1\le r<\omega.
\]
Every nonempty final interval of $\lambda$ has order type $\lambda$.
Every bounded initial interval of $\lambda$ has order type less than
$\lambda$. It follows that an increasing embedding of a word of length
$\lambda$ into another word of that length has cofinal image.

\subsection{The equal-block principle}

\begin{lemma}[Equal blocks force coordinate preservation]\label{lem:equalblocks}
Let $u_1,\ldots,u_d$ and $v_1,\ldots,v_d$ be words, all of length exactly
$\lambda$. Then
\begin{equation}\label{eq:equalblocks}
  u_1\cdots u_d\emb v_1\cdots v_d
  \quad\Longleftrightarrow\quad
  u_i\emb v_i\text{ for every }i.
\end{equation}
More generally, an embedding from $d$ such blocks into $e$ such blocks
requires $d\le e$.
\end{lemma}

\begin{proof}
For each source block, consider the last target block it meets. A tail of
the source block lies entirely in that target block. Its length is still
$\lambda$, so its image is cofinal there. Successive source blocks must
therefore have strictly increasing last target blocks: a later source
position cannot be placed after a cofinal image while remaining in the
same target block. This proves $d\le e$.

When $d=e$, the strictly increasing assignment from source block indices
to their last target block indices must be the identity. The first source
block is contained in the first target block. After its cofinal image,
no position of the second source block can lie in the first target block;
and its last target block is the second. Continue inductively. Thus each
source block lies wholly in its corresponding target block, proving
necessity in \eqref{eq:equalblocks}. Sufficiency is componentwise
concatenation of embeddings.
\end{proof}

Indecomposability of the individual words is not assumed in this lemma.
Their common ordinal length supplies the necessary cofinality.
Equal numbers of blocks, however, are essential to its coordinate conclusion.

\subsection{The protection condition}

\begin{definition}
A word is \emph{$\lambda$-ended} if its length is $\lambda d$ for some
positive finite integer $d$.
\end{definition}

Every nonempty tail of a $\lambda$-ended word has order type at least
$\lambda$. More precisely, its length is $\lambda e$ for some positive
finite $e$. A tail may contain several remaining $\lambda$-blocks; it
need not have length exactly $\lambda$.

\begin{lemma}[No bounded spill into a protected block]\label{lem:nospill}
Let $u$ be $\lambda$-ended, and let $z$ have length $\lambda$.
In an embedding, a nonempty tail of $u$ cannot be mapped into a bounded
initial interval of a target occurrence of $z$.
\end{lemma}

\begin{proof}
Such a tail has length at least $\lambda$, whereas the target interval
has length less than $\lambda$. This violates
\eqref{eq:lengthmonotone}.
\end{proof}

This tiny observation is the protection mechanism. Without it, an old
finite suffix can disappear into the beginning of a new infinite
separator, destroying coordinate reflection.

\section{Activating the first atom at a new level}\label{sec:activation}

The first atom of a new height needs a different construction from the
later atoms. At this point there is not yet a large family of
$\lambda$-ended words to place between separators.

Choose a minimal letter $a\in P$ and another letter $b\ne a$.
Then $b\not\le_P a$. Put
\[
  c=a^\lambda.
\]
This is an indecomposable of length $\lambda$. It is minimal among the
indecomposables of that length: any word of length $\lambda$ embedding
into it can use only letters at most $a$, hence only the letter $a$.

\begin{lemma}[A terminal marker protects the whole lower-level word]
\label{lem:marker}
The map
\begin{equation}\label{eq:marker}
  F\colon\W_r(P)\longrightarrow\{\text{words of length }\lambda\},
  \qquad F(u)=u\,b\,a^\lambda
\end{equation}
is an order embedding on equivalence classes.
Its image is contained in $\Lcal(\I_r(P)\cup\{c\})$.
\end{lemma}

\begin{proof}
Since $|u|<\lambda$ and $\lambda$ is an infinite additively indecomposable
ordinal, $|u|+1+\lambda=\lambda$. Concatenation proves monotonicity.

Suppose $u b a^\lambda\emb v b a^\lambda$. The distinguished source
occurrence of $b$ cannot map into the target $a^\lambda$ tail, because
$b\not\le_P a$. It must map either to the distinguished target $b$ or
into $v$. Every position of $u$ maps strictly before that position, and
therefore inside $v$. Hence $u\emb v$.

Finally, $u$ is a finite concatenation of lower-height atoms up to
equivalence, $b$ itself is a lower-height atom, and $a^\lambda=c$.
\end{proof}

It is not necessary that the distinguished source $b$ map to the
\emph{distinguished} target $b$. The proof deliberately allows it to land
earlier inside $v$; either possibility recovers $u\emb v$.

\begin{proposition}[First-atom amplification]\label{prop:activation}
Suppose $\otp(\W_r(P))=H_N$, where $N=N_r(P)$.
Then the language $\Lcal(\I_r(P)\cup\{a^\lambda\})$ contains a
$\lambda$-ended induced suborder of maximal order type $H_{N+1}$.
The language itself has the same maximal order type.
\end{proposition}

\begin{proof}
Let $D$ be the image of $F$. By \cref{lem:marker}, $\otp(D)=H_N$ and every
word in $D$ has length $\lambda$. For every $d\ge1$, concatenate $d$
independently chosen elements of $D$. By \cref{lem:equalblocks}, the
resulting induced suborder is a copy of $D^d$.

The images for different $d$ have distinct lengths $\lambda d$, so their
quotient classes are disjoint. Length comparison prohibits an embedding
from a higher-$d$ layer to a lower-$d$ layer. Their union is
$\lambda$-ended, and \cref{lem:layers} gives type at least $H_{N+1}$.
There are exactly $N+1$ available indecomposable generators, so
\cref{cor:upper} gives the matching upper bound for the whole language
and hence for this induced suborder.
\end{proof}

This is the only point in the main induction where the availability of
a second inequivalent letter is required. For a singleton alphabet no
such $b$ exists, and the much smaller order type in
\cref{cor:exception} is the actual answer.

\section{Separator amplification}\label{sec:separator}

We now prove the main lower-bound lemma. It permits new generators that
dominate old ones, so it is not merely a construction with pairwise
incomparable separators.

\begin{lemma}[Protected separators]\label{lem:separator}
Let $B$ be a finite set of indecomposable word classes, each of length at
most $\lambda$. Let $z$ be an indecomposable word of length $\lambda$ such that
\begin{equation}\label{eq:separatorcondition}
  z\nemb b\qquad\text{for every }b\in B.
\end{equation}
Let $E\subseteq\Lcal(B)$ be an induced suborder consisting of
$\lambda$-ended words. For $m\ge0$, define
\begin{equation}\label{eq:separator-map}
  \Phi_m(u_0,\ldots,u_m)=u_0 z u_1 z\cdots z u_m,
  \qquad u_i\in E.
\end{equation}
Then $\Phi_m$ is an order embedding of $E^{m+1}$ into
$\Lcal(B\cup\{z\})$. Moreover,
\begin{equation}\label{eq:separatorcount}
  \Phi_m(\vec u)\emb\Phi_t(\vec v)\quad\Longrightarrow\quad m\le t.
\end{equation}
All outputs are $\lambda$-ended.
\end{lemma}

\begin{proof}
First, $z$ does not embed into \emph{any} word of $\Lcal(B)$.
Indeed, otherwise \cref{lem:onefactor} would embed it into one generator,
contrary to \eqref{eq:separatorcondition}.

\emph{Step 1: each source separator has a target separator as its last
component.} Consider an embedding of an output of $\Phi_m$ into an output
of $\Phi_t$. Regard the latter as a finite concatenation of its displayed
old words $v_j$ and displayed new separators $z$. Fix one displayed
source separator and take the last displayed target component it meets.
A tail of the source separator embeds into this component. Since the
separator is indecomposable, the whole separator embeds into that component.
The component cannot be an old word $v_j$. It must be a displayed target
separator.

\emph{Step 2: the matching of separators is strictly increasing.}
The tail from Step~1 has length $\lambda$, so its image in that target
separator is cofinal. No position of a subsequent source separator can
follow that cofinal image while remaining in the same target separator.
Thus the last target separators of consecutive source separators are
strictly increasing. There are $t$ target separators, proving
\eqref{eq:separatorcount}.

\emph{Step 3: with equal counts, the old coordinates cannot spill.}
Suppose $m=t$. The increasing assignment between two $m$-element index
sets is the identity. Therefore the $i$th source separator ends cofinally
in the $i$th target separator, for every $i$.

Consider an old source word $u_i$. If $i>0$, the cofinal image of the
preceding source separator prevents every position of $u_i$ from landing
in or before the preceding target separator. If $i<m$, all positions
of $u_i$ precede the image of the first position of the next source
separator. That first image lies in or before the corresponding next
target separator, because that separator is its last target component.

If any position of $u_i$ entered the next target separator, all subsequent
positions of $u_i$ would also lie there, strictly below the image of the
first position of the next source separator. This would embed a nonempty
tail of $u_i$ into a bounded initial interval of a length-$\lambda$ block.
It is impossible by \cref{lem:nospill}. Therefore $u_i$ remains wholly
inside $v_i$. The same argument handles $u_0$ without a left boundary;
$u_m$ has no right separator, but the previous separator already forces
it into the final target component $v_m$.

Hence $\Phi_m(\vec u)\emb\Phi_m(\vec v)$ implies
$u_i\emb v_i$ for all $i$. Conversely, coordinatewise embeddings together
with identity embeddings on $z$ concatenate to an embedding of the
outputs. This proves order reflection and preservation.

Finally, finite concatenation of words whose lengths are positive finite
multiples of $\lambda$, interleaved with words of length $\lambda$, again
has length a positive finite multiple of $\lambda$.
\end{proof}

\begin{remark}[Why the protection assumption is indispensable]
Take the old language to be finite words and $z=a^\omega$. The source
coordinate $a$ does not embed into the empty target coordinate, yet
$a z\equiv z$. Thus merely requiring $z$ not to embed into an old word
is insufficient for coordinate reflection. It controls how many
separators there can be, but not absorption before a separator.
The $\lambda$-ended condition is exactly what rules out this failure.
\end{remark}

\begin{proposition}[One additional atom, one additional exponent step]
\label{prop:amplification}
Under the hypotheses of \cref{lem:separator}, suppose $|B|=M$ and
$\otp(E)=H_M$. Then $\Lcal(B\cup\{z\})$ contains a $\lambda$-ended induced
suborder of maximal order type $H_{M+1}$, and its own maximal order type
is $H_{M+1}$.
\end{proposition}

\begin{proof}
Take the union of the images of $\Phi_m$ for all $m\ge0$.
By \cref{lem:separator}, its layers are copies of $E^{m+1}$; an embedding
cannot decrease $m$. Different layers cannot be equivalent, since
applying the separator-count inequality in both directions would force
their indices equal. \Cref{lem:layers} therefore gives lower bound
$H_{M+1}$.

Condition \eqref{eq:separatorcondition} ensures that $z$ is not equivalent
to a generator in $B$, so the enlarged generator alphabet has $M+1$
elements. Apply \cref{cor:upper} for the matching upper bound. The union
is $\lambda$-ended by the last assertion of \cref{lem:separator}.
\end{proof}

\section{Proof of finite-alphabet exactness}\label{sec:mainproof}

We now assemble the argument for \cref{thm:main}.
The finiteness, canonical order, and recurrence were proved in
\cref{sec:atoms}; only equality in the order-type bound remains.

\begin{proof}[Proof of \cref{thm:main}]
Induct on $k$. For $k=1$, $\W_1(P)=P^*$ and $N_1(P)=n$, so the result
is the finite-alphabet Higman formula.

Assume
\[
  \otp(\W_r(P))=H_{N_r(P)}
\]
for some $r\ge1$, and set $\lambda=\omega^r$ and $N=N_r(P)$.
Let $Z$ be the finite poset of the new height-$r$ atoms:
\[
  Z=\I_{r+1}(P)\setminus\I_r(P).
\]
Choose a minimal letter $a\in P$. The constant word $a^\lambda$ is a
minimal element of $Z$, as shown in \cref{sec:activation}. Choose a linear
extension
\[
  z_1,z_2,\ldots,z_d
\]
of $Z$ that starts with $z_1=a^\lambda$.
Define
\[
  B_i=\I_r(P)\cup\{z_1,\ldots,z_i\}\qquad(1\le i\le d).
\]

By \cref{prop:activation}, $\Lcal(B_1)$ contains a $\lambda$-ended induced
suborder $E_1$ of type $H_{N+1}$. Suppose, for $i\ge2$, that
$\Lcal(B_{i-1})$ contains a $\lambda$-ended induced suborder $E_{i-1}$
of type $H_{N+i-1}$.

The atom $z_i$ cannot embed into a lower-height atom: its ordinal length
is larger. It cannot embed into any $z_j$ with $j<i$ either, because
$z_1,\ldots,z_d$ is a linear extension of the atom order. Hence
$z_i\nemb b$ for every $b\in B_{i-1}$.
Apply \cref{prop:amplification} to obtain a $\lambda$-ended induced
suborder $E_i\subseteq\Lcal(B_i)$ of type $H_{N+i}$.

After all $d$ new atoms have been added,
\[
  B_d=\I_{r+1}(P),\qquad N+d=N_{r+1}(P).
\]
Its concatenation language is all of $\W_{r+1}(P)$ up to equivalence,
by \cref{thm:normalform}. The constructed lower bound is
$H_{N_{r+1}(P)}$, and \cref{cor:upper} gives exactly the same upper bound.
This completes the induction.
\end{proof}

\subsection{The stronger induction invariant}

At every positive height the proof establishes slightly more than the
statement of the theorem. Once the first atom of that height is activated,
each intermediate language already realizes its full Higman upper bound
on words ending in an entire block of that height. This invariant is
what makes the next separator admissible without adding a new letter to
$P$.

The proof does \emph{not} exhibit an order embedding of the entire free
word order $\I_k(P)^*$ into $\W_k(P)$. It exhibits a sequence of induced
product layers large enough to attain the same maximal order type.
Equality of maximal order types is weaker than an isomorphism of these
orders, and absorption shows why that distinction is necessary.

\section{The binary cases, worked out explicitly}\label{sec:binary}

\subsection{Two incomparable letters}

Let $P=A_2=\{a,b\}$. At heights zero and one the atoms are
\[
  a,\quad b,\quad A=a^\omega,\quad B=b^\omega,\quad C=(ab)^\omega.
\]
The nontrivial relations are generated by
\[
  a<A<C,\qquad b<B<C.
\]
The two middle atoms $A,B$ are incomparable, as are the two letters.
There are exactly five distinct atoms, so the upper bound is $H_5$.

The lower-bound construction can be read as three concrete stages.
Start with finite binary words, whose type is $H_2=\omega^\omega$.
The map
\[
  u\longmapsto u b a^\omega
\]
embeds them into words of length $\omega$. Finite products of these blocks,
layered by the number of blocks, give a family $E_A$ of $\omega$-ended
words with type
\[
  \sum_{d\ge1}(\omega^\omega)^{\otimes d}
  =\omega^{\omega^2}=H_3.
\]
Only $a,b,A$ are needed as generators.

Next add $B$. It embeds into neither a finite letter nor $A$, so it does
not embed into any old word. Interleave $m$ copies of $B$ among $m+1$
independently chosen members of $E_A$. The protected-separator lemma gives
a new $\omega$-ended family $E_{A,B}$ of type
\[
  \sum_{m\ge0}H_3^{\otimes(m+1)}=H_4=\omega^{\omega^3}.
\]
Finally add $C$. It does not embed into any of $a,b,A,B$: an infinite
repetition containing both letters cannot embed into a monochromatic
infinite block. The same construction yields
\[
  \sum_{m\ge0}H_4^{\otimes(m+1)}=H_5=\omega^{\omega^4}.
\]
This matches the five-generator upper bound and proves the binary
antichain instance of \cref{cor:first}.

Notice that $C$ \emph{dominates} the old infinite atoms. This does not
invalidate its use as a separator. What matters is that $C$ does not
embed \emph{into} an old atom; protection prevents the old coordinates
from being swallowed by a bounded part of a target $C$.

\subsection{Two comparable letters}

For $C_2=\{0<1\}$, the height-zero and height-one atoms are
\[
  0,\quad1,\quad A=0^\omega,\quad B=1^\omega.
\]
Here $(01)^\omega\equiv1^\omega$, because every source letter is at most
$1$ and the $1$ positions occur cofinally. The atom order is a diamond:
$0$ is below the incomparable atoms $1,A$, both of which are below $B$.

The same activation map $u\mapsto u1\,0^\omega$ gives type $H_3$ after
$A$ is added. Adding $B$ as a protected separator gives $H_4$.
There are only four atoms, so
\[
  \otp(\W_2(C_2))=\omega^{\omega^3}.
\]
Thus two alphabets with the same finite maximal order type $2$ already
have different exact transfinite-word types.

\subsection{All finite alphabets at height two}

The height-one atoms correspond to nonempty order ideals of $P$.
An antichain $S\subseteq P$ represents the periodic word
$(\prod_{s\in S}s)^\omega$; its class depends precisely on the ideal
$\down S$. Therefore
\[
  N_2(P)=n+J(P)-1,
\]
which proves \cref{cor:first} from the main theorem.

\begin{corollary}[Extremal alphabets of a fixed finite size]
\label{cor:extremal}
For every $n$-element poset $P$, with $n\ge2$,
\begin{equation}\label{eq:extremal}
  \omega^{\omega^{2n-1}}
  \ \le\ \otp(\W_2(P))
  \ \le\ \omega^{\omega^{n+2^n-2}}.
\end{equation}
The lower endpoint is attained exactly by chains, and the upper endpoint
exactly by antichains.
\end{corollary}

\begin{proof}
Every linear extension of $P$ gives $n+1$ distinct prefix ideals, so
$J(P)\ge n+1$. If $P$ is not a chain, choose incomparable $x,y$.
The ideals $\down x$ and $\down y$ are incomparable, so at least one is
not in the chain of prefix ideals; hence $J(P)>n+1$.
A chain has precisely $n+1$ ideals.

There are at most $2^n$ subsets of $P$. All of them are ideals exactly
when $P$ is an antichain: a strict comparison would make the singleton
containing its larger element fail to be an ideal. Substitute these
bounds into \eqref{eq:first}. The finite exponent parameter is strictly
increasing, so the equality characterizations follow.
\end{proof}

\section{Effective enumeration and computed examples}\label{sec:computations}

\subsection{Canonical generation as a finite algorithm}

The recurrence \eqref{eq:count} comes with a direct algorithm. Store each
leaf by its label and each node by a sorted tuple of child identifiers.
At a new height, enumerate the antichains of the already constructed
poset; retain those whose largest child height is the preceding height;
and create the corresponding nodes. Compute comparisons by
\eqref{eq:leafleaf}--\eqref{eq:nodenode}, memoizing each ordered pair of
identifiers.

This is a decision procedure for all bounded-height atom comparisons.
It also gives an exact ordinal notation for the maximal order type:
a finite integer $N_k(P)$ together with the fixed expression
$\omega^{\omega^{N_k(P)-1}}$. No enumeration of actual infinite words is
required.

The finite objects can nevertheless become large. With $N$ existing
nodes, there are at most $N^2$ ordered pairs to memoize and at most $N^2$
child-pair checks per pair, giving a simple $O(N^4)$ arithmetic-operation
bound for the complete comparison matrix. Antichain enumeration can have
exponentially many outputs. The code uses integer bit masks and explicit
size caps; it makes no claim of a polynomial-time algorithm in $k$ or
in the original alphabet size.

\subsection{Two independent ideal counters}

The package counts $J(R)$ in two ways. Let $M$ be the set of remaining
vertices in an induced subposet.

The first method counts antichains as independent sets of the
comparability graph. For a chosen $x\in M$, an antichain either omits
$x$ or contains $x$ and omits every vertex comparable with $x$:
\begin{equation}\label{eq:antichaincount}
 A(M)=A(M\setminus\{x\})+
 A\bigl(M\setminus\{y\in M:y\le x\text{ or }x\le y\}\bigr).
\end{equation}
The second method directly partitions ideals according to whether they
contain $x$:
\begin{equation}\label{eq:idealcount}
 J(M)=J(M\setminus\up_M x)+J(M\setminus\down_M x).
\end{equation}
In the first branch, excluding $x$ forces exclusion of everything above
it. In the second, the entire downset of $x$ is fixed as present, and
one chooses an ideal in the remainder. Both recursions have value one
on the empty poset and are memoized. For small posets the code additionally
checks all subsets directly against the definition of an ideal.

The two counters share the mathematically specified atom order, but they
do not share an antichain-enumeration implementation. Their agreement is
therefore a meaningful check of the numerical computations, not an
independent proof of the atom-comparison theorem.

\subsection{Counts and exact ordinal parameters}

Write $A_n$ for the $n$-element antichain and $C_n$ for the $n$-element chain.
The following are exact values of $N_k(P)$; a dash means not computed in
this package.

\begin{table}[ht]
\centering
\renewcommand{\arraystretch}{1.2}
\begin{tabular}{@{}lrrrrrr@{}}
\toprule
Alphabet & $k=1$ & $k=2$ & $k=3$ & $k=4$ & $k=5$ & $k=6$\\
\midrule
$A_2$ & 2 & 5 & 11 & 28 & 136 & ---\\
$C_2$ & 2 & 4 & 7 & 12 & 22 & 50\\
$A_3$ & 3 & 10 & 43 & 1962 & --- & ---\\
$C_3$ & 3 & 6 & 12 & 28 & --- & ---\\
\bottomrule
\end{tabular}
\caption{Canonical indecomposable counts. Each entry $N$ gives
$\otp(\W_k(P))=\omega^{\omega^{N-1}}$.}
\label{tab:counts}
\end{table}

For example, the theorem and the table give
\[
\begin{aligned}
  \otp(\W_3(A_2))&=\omega^{\omega^{10}}, &
  \otp(\W_4(A_2))&=\omega^{\omega^{27}},\\
  \otp(\W_3(C_2))&=\omega^{\omega^6}, &
  \otp(\W_4(A_3))&=\omega^{\omega^{1961}}.
\end{aligned}
\]
The last reported value for each family is obtained from the two exact
ideal counts at the preceding level. Its full next-level atom poset is
not materialized. The CSV records this distinction explicitly.

The first binary counts can be verified without software. The five-element
$\I_2(A_2)$ has nine ideals not containing $C$: choose a prefix in each of
the two disjoint two-element chains $a<A$ and $b<B$. There is one more
ideal containing $C$, namely the whole poset. Thus $J(\I_2(A_2))=10$
and $N_3(A_2)=2+10-1=11$.

The six height-two atoms in this case can be represented by
\[
\begin{gathered}
  \Omega(\{A\}),\quad \Omega(\{B\}),\quad \Omega(\{C\}),\\
  \Omega(\{A,B\}),\quad \Omega(\{A,b\}),\quad \Omega(\{a,B\}).
\end{gathered}
\]
For the two-element chain, $\I_2(C_2)$ is the four-element diamond and
has six ideals, giving $N_3(C_2)=2+6-1=7$.

\section{Growth, interpretation, and limits of the result}\label{sec:scope}

\subsection{A direct iterated-exponential lower bound}

The exact formula also supplies a simple finite-alphabet growth result.
For $n\ge2$, define integers
\[
  w_0(n)=n,\qquad
  w_{r+1}(n)=\binom{w_r(n)}{\lfloor w_r(n)/2\rfloor}.
\]

\begin{proposition}\label{prop:width}
The height-$r$ atoms over $A_n$ contain an antichain of size $w_r(n)$.
Consequently, for $k\ge2$,
\begin{equation}\label{eq:widthbound}
 N_k(A_n)\ge n+2^{w_{k-2}(n)}-1,
 \qquad
 \otp(\W_k(A_n))\ge
 \omega^{\omega^{n+2^{w_{k-2}(n)}-2}}.
\end{equation}
\end{proposition}

\begin{proof}
The $n$ leaves give the height-zero antichain. Given an antichain $D$ of
height-$r$ atoms, the nodes $\Omega(S)$ for nonempty $S\subseteq D$ are
ordered exactly by subset inclusion: no distinct members of $D$ compare.
Choosing all subsets of size $\lfloor|D|/2\rfloor$ produces an antichain
of the claimed size at height $r+1$. For the stated $n$, this size is
positive and the subsets used are nonempty.

Also, all $2^{|D|}-1$ nonempty subsets give distinct nodes at the next
height. Use a height-$(k-2)$ antichain and retain those nodes together with
the $n$ leaves. This proves the first inequality in \eqref{eq:widthbound};
the second follows from \cref{thm:main}.
\end{proof}

The elementary bounds
\[
  \frac{2^m}{m+1}\le\binom{m}{\lfloor m/2\rfloor}\le2^m
\]
follow because the largest binomial coefficient is at least the average
of the $m+1$ coefficients. Thus, for fixed $k$ and growing $n$, the
finite exponent parameter has the expected iterated-exponential growth.
The upper estimate
\[
  N_{k+1}(P)\le n+2^{N_k(P)}-1
\]
follows from \eqref{eq:count}. These are statements about finite alphabets
and growth in their size; no sharp asymptotic equivalent is claimed.

\subsection{The quotient, not the literal set of words}

There may be uncountably many literal words of length $\omega$ over a
finite alphabet, but at each fixed finite level there are only countably
many equivalence classes of all words. Indeed, every class has a finite
code over the finite atom alphabet. The large ordinals in the theorem
measure ordering complexity of this quotient, not the cardinality of
the original set of functions.

The integer $N_k(P)$ should also not be mistaken for the number of all
word classes. It counts only indecomposable classes. Their finite
concatenations have infinitely many classes and supply the much larger
maximal order type.

\subsection{Why this does not compute the all-height union}

Although
\[
  \W_{<\omega^\omega}(P):=\{u:|u|<\omega^\omega\}
  =\bigcup_{k\ge1}\W_k(P),
\]
one cannot infer its maximal order type merely by taking the supremum
of the numbers in \cref{thm:main}. Maximal order type is not generally
continuous under increasing unions of induced suborders.

A simple example is the increasing union of the finite squares
$\{0,\ldots,m-1\}^2$ under componentwise order. Each has maximal order
type $m^2$, so the supremum of their types is $\omega$. The union is
$\omega\times\omega$, whose maximal order type is
$\omega\otimes\omega=\omega^2$. The all-height problem requires its own
argument; the broader tree correspondence and all-height results in
\cite{ChopraPakhomov} are separate from the finite-level proof here.

\subsection{What remains outside the theorem}

The proof uses finiteness in two different ways. The normal-form step
needs only finitely many lower-height atom classes, so that recurrent
types can be isolated after a finite prefix. The exactness step adds all
new atoms in a finite linear extension and uses the finite-letter Higman
formula. Neither argument directly provides an exact answer for every
infinite wqo alphabet, arbitrary transfinite height, or arbitrary cutoff
ordinal that is not a power $\omega^k$.

Nor does a large maximal order type automatically specify an optimal
length-function bound for finite controlled bad sequences. Such a bound
also requires a norm and a control function. The ordinal theorem is an
input to that separate analysis, not a substitute for it.

\section{Verification, proof audit, and reproducibility}\label{sec:audit}

\subsection{What was actually checked}

The code was run successfully under Python~3.13.5, using only the standard
library. It is written to support Python~3.9 or later. The fixed random
seed is \texttt{20260919}. The recorded run reports:

\begin{table}[ht]
\centering
\renewcommand{\arraystretch}{1.13}
\begin{tabular}{@{}lr@{}}
\toprule
Check & Count\\
\midrule
All labelled posets on at most three elements & 24\\
Canonical atom posets through $k=3$ for those alphabets & 72\\
Agreements of two independent ideal counters & 87\\
Agreements with direct exhaustive subset counting & 64\\
Marker-map comparison pairs & 1,922\\
Equal-block comparison pairs & 10,000\\
Separator comparison pairs & 18,000\\
\quad of which equal-separator coordinate checks & 13,489\\
Singleton length-order comparison pairs & 625\\
Absorption sanity checks & 2\\
Exact count-table entries & 19\\
\bottomrule
\end{tabular}
\caption{The completed reproducibility run. These are finite checks, not
a formal verification of the transfinite proof.}
\label{tab:verification}
\end{table}

The count ``24'' includes one empty poset, one one-element poset, three
labelled two-element posets, and nineteen labelled three-element posets.
The atom checks verify reflexivity, antisymmetry, transitivity,
height monotonicity, and antichain reduction. The symbolic word checks
operate on exact ultimately periodic $\omega$-blocks; they do not replace
$\omega$ by a large finite integer. \Cref{app:decision} describes the
embedding algorithm and its correctness argument.

These checks exercise the finite representations, detect arithmetic and
implementation mistakes, and test many instances of the structural
lemmas. They do not establish arbitrary-height statements by exhaustive
search. The latter are supported by the written induction and cofinality
proofs.

\subsection{The main points audited in the proof}

The first possible failure is assuming that a surjective concatenation
map is order-reflecting. It need not be; the proof uses it only for the
upper bound. Order reflection for the lower bound is proved separately
on the protected families.

The second is confusing ``has the same cardinality'' with ``has the same
ordinal length.'' All no-spill arguments compare ordinal order types.
A tail of a $\lambda$-ended word has length at least $\lambda$, and cannot
fit below a single position of a target $\lambda$-block.

The third is silently treating an indecomposable as wholly contained in
its last target factor. The original embedding may straddle factors.
\Cref{lem:onefactor} extracts an equivalent tail to obtain a new
embedding; \cref{lem:separator} then uses cofinality of that tail in its
last target separator.

The fourth is choosing separators in the wrong order. The requirement
is that the \emph{new} atom not embed into any \emph{old} atom. An
increasing linear extension of the same-height atom poset supplies
exactly that direction. Old atoms are allowed to embed into the new one.

The fifth is combining product layers without checking cross-layer
comparisons. In activation, length makes the layer index monotone. In
separator amplification, \cref{eq:separatorcount} does so. This also
precludes identifying outputs from two different layers on the quotient.

Finally, the singleton alphabet is checked separately. It has one atom
at each height, so $N_k(1)=k$, but its word type is $\omega^k$, not
$H_k$ for $k\ge2$. The missing second letter is a genuine mathematical
obstruction to the activation argument, not a cosmetic exception.

\subsection{How to reproduce the package}

From the archive root, run
\begin{lstlisting}
python code/verify.py --out data
\end{lstlisting}
The program regenerates \texttt{data/verification.json}, the CSV count
table, and explicit JSON atom posets. To inspect a particular alphabet,
for example the binary antichain through height two, run
\begin{lstlisting}
python code/atoms.py --kind antichain --size 2 --k 3 --output atoms.json
\end{lstlisting}
A custom finite poset may be supplied as JSON containing distinct
\texttt{labels} and a full reflexive Boolean \texttt{relation} matrix.
The default cap prevents materialization of more than 5,000 atoms; the
cap can be changed explicitly. Counting can still be expensive before
a large object is materialized.

Build this document using
\begin{lstlisting}
latexmk -pdf -interaction=nonstopmode -halt-on-error article.tex
\end{lstlisting}
The bibliography is included in the source, so no bibliography database
or network access is needed. A separate \texttt{references.bib} is
provided for reuse. The original third-party papers and font files are
not included.

\appendix
\section{The finite-alphabet Higman formula}\label{app:higman}

This appendix records a proof of the finite case of
\eqref{eq:higman}, using Higman's wqo theorem and the standard calculus
\eqref{eq:wqocalculus}. It also explains why the separator construction
in the main argument has the right ordinal shape.

\begin{theorem}[Classical finite-alphabet formula]
For a nonempty finite poset $R$ with $n$ elements,
\[
  \otp(R^*)=\omega^{\omega^{n-1}}=H_n.
\]
\end{theorem}

\begin{proof}
Induct on $n$. For $n=1$, words are ordered by their finite lengths, so
the type is $\omega=H_1$.
Assume $n\ge2$ and the result for smaller alphabets.

\emph{Lower bound.} Choose a maximal letter $b\in R$ and set
$Q=R\setminus\{b\}$. In words
\[
  u_0 b u_1 b\cdots b u_m,\qquad u_i\in Q^*,
\]
every source $b$ must map to a target $b$, since $b\not\le q$ for
$q\in Q$. The number of displayed $b$'s cannot decrease under embedding.
When the two words have the same number, all displayed $b$'s match in
order, forcing the intervening coordinates into their corresponding
intervals. Consequently, the layer with $m$ occurrences of $b$ is an
induced copy of $(Q^*)^{m+1}$. The layers are disjoint and their index is
monotone, so \cref{lem:layers}, with $\otp(Q^*)=H_{n-1}$, gives lower
bound $H_n$.

\emph{Upper bound.} Fix a forbidden nonempty word
$w=p_1\cdots p_\ell$. A word that does not contain $w$ as an embedded
subsequence has a greedy factorization
\[
  u_1 b_1 u_2 b_2\cdots u_j b_j u_{j+1},\qquad 0\le j<\ell,
\]
where $b_i\ge p_i$ and $u_i$ uses only letters in
$Q_i=R\setminus\up p_i$. Each $Q_i$ is a proper subalphabet.
This factorization is obtained by greedily matching $p_1$, then $p_2$,
and so on, stopping when the next letter cannot be matched.

For a fixed $j$ and fixed matched letters $b_i$, concatenation is a
monotone surjection from a product of $j+1$ orders $Q_i^*$ onto this
part of the residual. By induction every nonempty $Q_i^*$ has type at
most $H_{n-1}$; an empty $Q_i$ contributes the one-element order,
which also satisfies the bound. There are only finitely many choices
of $j$ and of the $b_i$. Thus the residual
$R^*\setminus\up w$ is a monotone image of a finite disjoint union of
products, whose maximal order type is bounded by a finite natural sum
of ordinals at most
\[
  H_{n-1}^{\otimes\ell}
  =\omega^{\omega^{n-2}\cdot\ell}.
\]
Every such finite sum is strictly less than $H_n$.
The residual for the empty forbidden word is empty.
Using the residual identity in \eqref{eq:wqocalculus} and the fact that
$H_n$ is a limit ordinal gives $\otp(R^*)\le H_n$.
\end{proof}

The transfinite proof differs in one crucial respect. A new separator
is an infinite word, not a single letter. Old coordinates can in
principle enter its beginning even when its occurrences have all been
matched. Protected tails are what restore the coordinate conclusion.

\section{Exact symbolic embedding below \texorpdfstring{$\omega^2$}{omega squared}}\label{app:decision}

The test program represents a word as
\begin{equation}\label{eq:symbolic}
  p_1 c_1^\omega\ p_2 c_2^\omega\ \cdots\ p_t c_t^\omega\ q,
\end{equation}
where each $p_i,q$ is a finite word and each $c_i$ is a finite nonempty
cycle. This is an exact symbolic representation of a transfinite word,
not a truncation. It is sufficiently general to represent every class
below $\omega^2$ over a finite alphabet, by the normal-form argument,
although the tests use bounded finite parameter sets.

\subsection{The greedy algorithm}

Maintain a cursor in the target. Within a target block it consists of
its block index and a finite offset in its finite prefix or periodic
tail. After the final infinite block it is an offset in the finite suffix.
Process a source finite prefix one letter at a time, always choosing
the earliest target position carrying a label above the source label.
In a periodic tail, checking one cycle suffices to find the earliest
match; if no cycle letter is suitable, skip to the next target block.

For a source periodic tail $c^\omega$, find the earliest target infinite
block, at or after the cursor, whose recurring cycle $d$ satisfies
\begin{equation}\label{eq:cyclecondition}
  (\forall a\text{ occurring in }c)
  (\exists b\text{ occurring in }d)\ a\le_P b.
\end{equation}
If there is no such block, embedding fails. Otherwise embed the periodic
tail cofinally into that block and advance the cursor to the beginning
of the next target block. Continue with the next source prefix. After
all source infinite blocks are consumed, greedily process the finite
source suffix.

\subsection{Why the algorithm is exact}

Greedy matching of a finite letter is safe because replacing its target
position by an earlier suitable one can only enlarge the remaining
target interval. This is the usual finite-subsequence argument and
applies unchanged inside a transfinite target.

An infinite periodic source tail has infinitely many occurrences of
every label in its cycle. If it embeds into a finite concatenation of
target infinite blocks and a finite suffix, its image ends cofinally
in one target infinite block. A source tail of the same periodic type
therefore embeds into that block's recurring tail. This forces
\eqref{eq:cyclecondition}; finite exceptional target prefixes cannot
supply infinitely many occurrences of a missing label.

Conversely, \eqref{eq:cyclecondition} suffices: place successive source
letters in suitable positions of successive sufficiently late target
cycles. Any finite cursor offset is harmless. Choosing the earliest
compatible target block leaves the largest possible collection of
whole target blocks for subsequent source material. No subsequent
source position could remain in the same target block after a cofinal
image. Thus advancing to the next block is both sound and optimal.
Induction over the finitely many source blocks proves correctness.

\subsection{Test families and negative controls}

The marker tests compare all finite binary source words of length at
most four with all finite binary target words of length at most four,
for both the chain and antichain letter orders. There are
$2(1+2+4+8+16)^2=1{,}922$ pairs.

The equal-block tests use up to four exact $\omega$-blocks, with finite
prefixes and one of the cycles $0$, $1$, or $01$. The separator tests
first add $1^\omega$ to a language whose infinite blocks have cycle $0$.
For the antichain they also add $(01)^\omega$ to a language whose infinite
blocks are monochromatic. Both equal and unequal separator counts are
tested, including positive comparisons and self-comparisons.

As negative controls against an incorrectly rigid concatenation
implementation, the program verifies the genuine absorption
$a a^\omega\equiv a^\omega$. It also checks that over a singleton
alphabet, the embedding decision agrees exactly with comparison of the
ordinal lengths $\omega t+s$ on the finite test range.

\begin{thebibliography}{9}
\bibitem{Altman}
Harry Altman.
\emph{Bounding finite-image sequences of length $\omega^k$}.
arXiv:2409.03199v2, revised March 10, 2026.
\href{https://arxiv.org/abs/2409.03199v2}{arXiv:2409.03199v2}.
The specific target is Example~3.15, page~10.

\bibitem{ChopraPakhomov}
Alakh Dhruv Chopra and Fedor Pakhomov.
\emph{Well-quasi-orders on finite trees and transfinite sequences}.
arXiv:2602.09830v1, February 10, 2026.
\href{https://arxiv.org/abs/2602.09830v1}{arXiv:2602.09830v1}.
See especially Section~4.3 and Theorem~4.26.

\bibitem{deJonghParikh}
D. H. J. de Jongh and R. Parikh.
Well-partial orderings and hierarchies.
\emph{Indagationes Mathematicae (Proceedings)} 80(3) (1977), 195--207.
\href{https://doi.org/10.1016/1385-7258(77)90067-1}
{doi:10.1016/1385-7258(77)90067-1}.

\bibitem{Higman}
Graham Higman.
Ordering by divisibility in abstract algebras.
\emph{Proceedings of the London Mathematical Society} s3-2(1) (1952), 326--336.

\bibitem{Schmidt}
D. Schmidt.
Well-partial orderings and their maximal order types.
In P. M. Schuster, M. Seisenberger, and A. Weiermann, editors,
\emph{Well-Quasi Orders in Computation, Logic, Language and Reasoning},
Springer, 2020, pages 351--391. Reprint of the 1979 habilitation thesis.
\end{thebibliography}

\end{document}
